% TOP_PROBLEM: {"id": 30007064, "title": "Invariant anticanonical sections under structure-sheaf acyclicity", "category_id": 5, "category": {"id": 5, "name": "algebraic_geometry", "display_name": "Algebraic Geometry", "description": "Geometric objects defined by polynomial equations.", "slug": "algebraic-geometry", "order_index": 5, "created_at": "2026-07-31T15:26:25.671Z"}, "status": "open"}
% FIRST_POSTED: 2026-10-06
% ATTEMPT_STATUS: solved
% OPENAI_CLAIM: {"schema_version":1,"source_repo":"https://github.com/openai/math","source_commit":"adc7f1241b42e322a6451854ab7e4b4c146bf78a","release_date":"2026-10-06","match":"full","resolution":"proved","independently_reviewed":false,"second_pass":"agreed","families":[{"family":"068","title":"Anticanonical nonvanishing in every dimension","subject":"Algebraic and complex geometry","manuscripts":[{"title":"Exact orders and invariant anticanonical linear systems","pdf_path":"preprints/Exact-orders-and-invariant-anticanonical-linear-systems-September-26-2026/main.pdf","readme_path":"preprints/Exact-orders-and-invariant-anticanonical-linear-systems-September-26-2026/README.md","date":"2026-09-26","theorem_ref":"V:invariant-main (source/sections/introduction.tex:144)","statement_sources":[{"path":"preprints/Exact-orders-and-invariant-anticanonical-linear-systems-September-26-2026/build/source/sections/introduction.tex","line":144,"label":"V:invariant-main"}]}],"lean":null,"reasoning_trace":null}]}
\documentclass[11pt]{article}
\usepackage[margin=25mm]{geometry}
\usepackage{fontspec}
\setmainfont{Latin Modern Roman}
\usepackage{amsmath,amssymb,mathtools}
\usepackage{unicode-math}
\setmathfont{Latin Modern Math}
\usepackage{enumitem}
\usepackage{xurl,hyperref,bookmark}
\hypersetup{colorlinks=true,urlcolor=blue}
\setcounter{secnumdepth}{0}
\setlength{\parindent}{0pt}
\setlength{\parskip}{5pt}
\setlength{\emergencystretch}{3em}
\begin{document}
\section*{Invariant anticanonical sections under structure-sheaf acyclicity}

\subsection{OpenAI's claim}
Asymptotically zero curvature-null valuation for a smoothly semipositive anticanonical line is claimed to be attained exactly by an effective rational anticanonical divisor. The cited manuscript argues for that result; its relation to this problem is specified below.

\subsection{Mathematical statement of the claimed result}
\paragraph{Mathematical claim 1.} Let $X$ be smooth connected projective complex, with a smooth semipositive anticanonical metric and $H^i(X,\mathcal O_X)=0$ for every $i>0$. Given an algebraic torus action by $T$ on $X$, use its differential linearization on $-K_X$, OpenAI claims
\[
\exists m>0:\quad H^0(X,\mathcal O_X(-mK_X))^T\ne0.
\]
The section is invariant under the natural action, not merely an arbitrary character eigensection; smooth semipositivity and the cohomological vanishings are retained.

\subsection{Outline of the argument}
Manuscript 1 (family 068, 2026-09-26): In a finite bigraded ring, sections approaching a boundary ray use off-ray generators in negligible proportion. A boundary generator with suitably small normalized valuation is selected before the exceptional correction is pushed away; centeredness upgrades the resulting upper bound to exact order zero.

\subsection{Scope relative to this problem}
Manuscript 1 (family 068, 2026-09-26): The problem selects the mathematical statement summarized here, including its stated hypotheses and parameter range. The cited theorem gives a claimed proof of that assertion in full.

\subsection{Formal verification}
Family 068 has no Lean scope page in this snapshot; no formalization of these manuscript results is listed.

This site has not executed the Lean files or checked their dependencies. These descriptions report the scope of the linked declarations, rather than a new verification of the proofs.

\subsection{Status caveat}
OpenAI attributes these results to an internal model. Its release notice cautions that results lacking full formalization can contain mistakes. This site has not independently assessed the proofs; a solved label records the scope of the reported claim.

\subsection{Sources}
\begin{itemize}
\item OpenAI, \emph{Exact orders and invariant anticanonical linear systems}, 2026-09-26; citation key \texttt{OAI:Exact-orders-and-invariant-anticanonical-linear-systems-September-26-2026}. \href{https://github.com/openai/math/blob/adc7f1241b42e322a6451854ab7e4b4c146bf78a/preprints/Exact-orders-and-invariant-anticanonical-linear-systems-September-26-2026/main.pdf}{PDF (GitHub)}; \href{https://github.com/openai/math/raw/adc7f1241b42e322a6451854ab7e4b4c146bf78a/preprints/Exact-orders-and-invariant-anticanonical-linear-systems-September-26-2026/main.pdf}{PDF (direct)}; \href{https://github.com/openai/math/blob/adc7f1241b42e322a6451854ab7e4b4c146bf78a/preprints/Exact-orders-and-invariant-anticanonical-linear-systems-September-26-2026/README.md}{Manuscript page and citation}; \href{https://github.com/openai/math/blob/main/preprints/Exact-orders-and-invariant-anticanonical-linear-systems-September-26-2026/main.pdf}{Latest version}; Statement source: \href{https://github.com/openai/math/blob/adc7f1241b42e322a6451854ab7e4b4c146bf78a/preprints/Exact-orders-and-invariant-anticanonical-linear-systems-September-26-2026/build/source/sections/introduction.tex\#L144}{V:invariant-main}.
\item \href{https://github.com/openai/math/blob/adc7f1241b42e322a6451854ab7e4b4c146bf78a/CONTENTS.md}{OpenAI catalogue entry}.
\end{itemize}
\end{document}
